\subsection{Diseño del reconocimiento y de la evaluación mediante CLS}
\label{sec:vision-metricas}

La evaluación ortográfica se basa en CLS, presentada en la Sección~\ref{sec:cls}. El diseño separa tres responsabilidades: obtener una transcripción con Google Cloud Vision, compararla con las respuestas aceptadas del ejercicio y someter el resultado a revisión del tutor. El cálculo CLS es determinista y se ejecuta en el servidor; no requiere entrenar un modelo lingüístico ni interpretar el significado de la respuesta.

\subsubsection{Flujo de registro del intento}
\begin{enumerate}
    \item El alumno envía \texttt{id\_ejercicio\_tutor} e \texttt{imagen\_codificada} mediante \texttt{POST /api/alumno/intento}. La dependencia de autenticación identifica al alumno a partir del JWT.
    \item \texttt{registrar\_intento\_dao}, en \texttt{Modelo/dao\_alumno.py}, verifica la pertenencia de la asignación al tutor del alumno y su estado activo. En dictados también verifica la fecha de vencimiento.
    \item \texttt{Servicios/ocr.py} valida la imagen y, si el reconocimiento está habilitado, solicita \texttt{document\_text\_detection} a Google Cloud Vision. Recupera el texto y la anotación con la geometría de palabras y símbolos.
    \item \texttt{Servicios/evaluacion.py} selecciona la referencia y llama a \texttt{calcular\_cls}, definido en \texttt{Servicios/cls.py}. En paralelo conceptual, la función de legibilidad usa la anotación y la imagen para proponer medidas del trazo; ambas operaciones se ejecutan dentro de la misma solicitud.
    \item El DAO inserta en \texttt{Intentos} la imagen, la transcripción y \texttt{resultado\_ocr\_json}. Este JSON conserva \texttt{evaluacion\_preliminar}, incluido su objeto \texttt{cls}, y \texttt{legibilidad}. La respuesta HTTP 201 contiene el mensaje de registro y el objeto \texttt{ocr}.
\end{enumerate}

El reconocimiento y los cálculos preceden a la inserción del intento. No hay una cola de evaluación en segundo plano: si el OCR genera una excepción, la entrega no se inserta. Cuando está desactivado o devuelve una respuesta válida sin texto, se permite guardar el intento para revisión manual. Registrar una entrega no crea por sí mismo un análisis confirmado ni modifica los promedios.

\subsubsection{Preparación de imagen y contrato del OCR}
\texttt{preparar\_imagen} acepta PNG o JPEG como URL de datos en Base64. Aplica límites de 10 MB a los bytes de entrada y al PNG procesado y de 20 millones de píxeles a la imagen. Compone las transparencias sobre blanco y genera un PNG RGB con las dimensiones de entrada.

La variable \texttt{GOOGLE\_VISION\_ENABLED} habilita la llamada externa. El cliente utiliza credenciales predeterminadas de aplicación disponibles en el servidor, un tiempo de espera de 30 segundos y \texttt{retry=None}. \texttt{GOOGLE\_VISION\_LANGUAGE\_HINTS} permite configurar pistas de idioma: el código usa \texttt{es} cuando no hay un valor y una lista vacía cuando se establece \texttt{auto}. Estas opciones no corrigen la escritura del alumno.

El servicio devuelve \texttt{proveedor}, \texttt{estado}, \texttt{texto} y \texttt{full\_text\_annotation}. Los estados son \texttt{desactivado}, \texttt{procesado} y \texttt{sin\_texto}. La anotación conserva la jerarquía de páginas, bloques, párrafos, palabras y símbolos. Sus cajas delimitadoras son información del reconocimiento, no máscaras exactas de la tinta.

Los errores de imagen se comunican con 422 o 413; los problemas de credenciales o disponibilidad, con 503; y una respuesta del proveedor con un mensaje de error, con 502. El navegador no recibe las credenciales de Google.

\subsubsection{Selección de la respuesta de referencia}
La configuración se obtiene de \texttt{contenido\_base.evaluacion}. Los modos admitidos son \texttt{copia}, \texttt{completar}, \texttt{imagen} y \texttt{dictado}, con una lista no vacía de \texttt{respuestas\_aceptadas}. La comparación normaliza Unicode a NFC, trata los guiones entre caracteres de palabra como separadores y separa palabras y signos. Cuando hay varias respuestas, elige la de mayor similitud de tokens mediante \texttt{SequenceMatcher}, con \texttt{autojunk=False}; en un empate conserva la primera. Esta selección se realiza antes del conteo CLS y no equivale a elegir el mayor puntaje CLS.

La comparación también produce observaciones de omisión, adición, sustitución, acentuación, mayúsculas y puntuación. Son indicios para contrastar con la imagen. El número de símbolos con confianza menor que 0.8 ayuda a orientar la revisión, sin expresar una probabilidad de error del alumno.

\subsubsection{Reglas del cálculo CLS implementado}
\label{sec:diseno-cls}
El algoritmo extrae secuencias de letras después de normalizar a NFC. Conserva la distinción entre mayúsculas, acentos y la letra ñ, y excluye números y signos de puntuación. Del manual de \textit{aimswebPlus} se toma el criterio de contar los pares adyacentes e incluir los límites de cada palabra \cite[pp.~62--63]{aimswebplus2025manual}. Para una palabra esperada $w$, hay $|w|-1$ pares interiores y dos pares en los extremos, por lo que el total es $|w|+1$. Por ejemplo, \texttt{casa} aporta los pares inicio--c, c--a, a--s, s--a y a--fin.

Para automatizar el criterio, Graf-IA registra las ocurrencias de cada par en dos multiconjuntos, $B(w)$ para la palabra esperada y $B(v)$ para la escrita. Cada coincidencia consume una ocurrencia en ambos, de modo que el número de veces que puede acreditarse un par $b$ es la menor de sus dos frecuencias. La Ecuación~\eqref{eq:cls-pares} suma esas coincidencias y describe el cálculo implementado en \texttt{Servicios/cls.py}.
\cite{aimswebplus2025manual}.
\begin{equation}
 C(w,v)=\sum_b \min\bigl(\operatorname{freq}_{B(w)}(b),
                            \operatorname{freq}_{B(v)}(b)\bigr).
 \label{eq:cls-pares}
\end{equation}
La coincidencia exige adyacencia y orden dentro del par. El conteo por frecuencias admite coincidencias en distintas posiciones de la palabra.

Las listas de palabras se alinean mediante \texttt{SequenceMatcher}. En los segmentos de sustitución se emparejan en orden. Una palabra esperada sin correspondencia aporta cero secuencias correctas y las palabras añadidas quedan fuera del cálculo. En la Ecuación~\eqref{eq:cls-implementado}, $C$ suma los aciertos y $P$ suma los pares posibles de las $m$ palabras esperadas. El cociente pondera cada palabra según su número de pares posibles y se convierte mediante la regla de la Ecuación~\eqref{eq:cls-escala}.
\cite{aimswebplus2025manual}.
\begin{equation}
 C=\sum_{i=1}^{m}C(w_i,v_i),\qquad
 P=\sum_{i=1}^{m}(|w_i|+1),\qquad
 S_{\mathrm{CLS}}=10\frac{C}{P}\quad(P>0).
 \label{eq:cls-implementado}
\end{equation}
El servicio redondea el puntaje sugerido a dos decimales y retorna \texttt{correctas}, \texttt{posibles}, \texttt{puntaje\_sugerido} y el detalle \texttt{por\_palabra}. Si no hay secuencias posibles, el puntaje es \texttt{null}. La Tabla~\ref{tab:ejemplos-cls-diseno} muestra ejemplos de transcripciones y sus puntajes CLS.

\begin{table}[H]
\centering
\caption{Ejemplos del comportamiento del cálculo CLS}
\label{tab:ejemplos-cls-diseno}
\begin{tabular}{llll}
\toprule
Referencia & Transcripción & $C/P$ & Puntaje \\
\midrule
casa & casa & $5/5$ & 10.00 \\
casa & caza & $3/5$ & 6.00 \\
casa & Casa & $3/5$ & 6.00 \\
canción & cancion & $6/8$ & 7.50 \\
perro & pero & $5/6$ & 8.33 \\
El gato. & El gato & $8/8$ & 10.00 \\
\bottomrule
\end{tabular}
\end{table}

Estas reglas delimitan el indicador: no mide puntuación, gramática, significado ni calidad del trazo. Una palabra añadida puede generar una observación sin reducir el CLS; por ejemplo, \texttt{El perro grande come} frente a \texttt{El perro come} obtiene $14/14$. Por ello, el puntaje y las observaciones se presentan como datos distintos y se conservan junto con la imagen original.

\subsubsection{Estados y revisión por el tutor}
La evaluación textual devuelve \texttt{no\_configurada} si falta una referencia válida, \texttt{sin\_texto} si no hay transcripción utilizable y \texttt{revision\_manual} si el texto o alguna referencia supera 10\,000 caracteres. Solo el estado \texttt{preliminar} incorpora el CLS calculado sobre la transcripción; todos los estados indican \texttt{requiere\_revision=true}.

Para consultar intentos anteriores, \texttt{cls\_de\_intento} reutiliza el CLS guardado o lo calcula al consultar si hay una referencia válida. Cuando no existe texto puede obtener un conteo de cero y el total posible de la primera referencia para habilitar la captura manual del tutor. Ese valor no confirma un desempeño de cero ni se incorpora automáticamente al progreso.

El tutor contrasta la imagen con la referencia y envía \texttt{cls\_correctas} mediante \texttt{PUT /api/tutor/calificar/\{id\_intento\}}. El servidor verifica la pertenencia del intento, obtiene las secuencias posibles desde la evaluación conservada o la referencia y exige $0\leq C\leq P$. El cliente no fija el denominador. El conteo confirmado actualiza o inserta \texttt{Analisis\_Ortografico}; su \texttt{ortografia\_score} se calcula con la Ecuación~\ref{eq:cls-implementado}.

\noindent\begin{minipage}{\linewidth}
\begin{lstlisting}[language=JSON, caption={Ejemplo de CLS confirmado en sugerencias\_json}, label={lst:ocr-output}]
{
  "metrica": "CLS",
  "correctas": 4,
  "posibles": 5,
  "sugerido": {
    "correctas": 3,
    "posibles": 5,
    "puntaje_sugerido": 6.0
  },
  "confirmado_por_tutor": true
}
\end{lstlisting}
\end{minipage}

Los conteos y la confirmación son claves de \texttt{sugerencias\_json}, no columnas independientes. \texttt{cantidad\_errores} almacena $P-C$, es decir, secuencias no acreditadas; no representa un número de palabras incorrectas. La calificación general y la retroalimentación se guardan por separado en \texttt{Intentos}. Omitir \texttt{cls\_correctas} no crea ni elimina una confirmación anterior.

\subsubsection{Legibilidad como evaluación complementaria}
\label{sec:preparacion-metricas}
\texttt{Servicios/legibilidad.py} ya propone cuatro medidas independientes del CLS: alineación, tamaño, espaciado e inclinación. Agrupa palabras en renglones por solapamiento vertical de sus cajas y utiliza las letras reconocidas como apoyo geométrico. La implementación no mide uniformidad del grosor.

\begin{itemize}
    \item \textbf{Alineación:} ajusta una recta por mínimos cuadrados a la base de letras sin descendentes, con al menos tres puntos por renglón. Calcula la raíz de la media de los residuos al cuadrado, la divide entre la altura mediana y combina los renglones ponderándolos por su número de puntos.
    \item \textbf{Tamaño:} calcula el coeficiente de variación de alturas de al menos tres minúsculas del conjunto de letras comparables definido en el servicio.
    \item \textbf{Espaciado:} calcula el coeficiente de variación de los huecos entre palabras; si hay menos de dos, utiliza los huecos entre letras. Si la media de los huecos es no positiva, propone cero.
    \item \textbf{Inclinación:} analiza píxeles con intensidad menor que 128 dentro de las cajas de letras y estima un ángulo respecto a la vertical mediante momentos de la tinta. Promedia los ángulos absolutos ponderados por la cantidad de tinta.
\end{itemize}

La conversión a un puntaje de legibilidad es una regla de elaboración propia, implementada en \texttt{Servicios/legibilidad.py}. Para cada desviación $d$, se fija un puntaje de 10 en el umbral $t_{\mathrm{bueno}}$ y uno de 0 en $t_{\mathrm{malo}}$, con $t_{\mathrm{bueno}}<t_{\mathrm{malo}}$. La interpolación lineal mide la distancia de $d$ al umbral de cero, $t_{\mathrm{malo}}-d$, como fracción del intervalo entre los dos umbrales. La Ecuación~\eqref{eq:legibilidad-lineal} limita esa fracción al intervalo de 0 a 1 y la multiplica por 10.
\begin{equation}
 Q(d)=10\max\left(0,\min\left(1,
 \frac{t_{\mathrm{malo}}-d}{t_{\mathrm{malo}}-t_{\mathrm{bueno}}}\right)\right).
 \label{eq:legibilidad-lineal}
\end{equation}
El servicio redondea la sugerencia a un decimal. El intervalo entre umbrales va de $0.05$ a $0.35$ en alineación, de $0.10$ a $0.50$ en tamaño, de $0.15$ a $0.80$ en espaciado y de $10$ a $40$ grados en inclinación. Estos parámetros del código están pendientes de calibración pedagógica.

Una medida sin información suficiente se representa con \texttt{null}. La sugerencia global promedia solo las medidas disponibles y conserva los valores crudos en \texttt{detalle}. Para confirmar legibilidad, el contrato exige las cuatro medidas entre 0 y 10; el DAO las almacena en \texttt{Analisis\_Caligrafico}, con el contexto sugerido en \texttt{observaciones}. El OCR, CLS y la legibilidad deben validarse por separado según la Sección~\ref{sec:validacion-vision}.
